\section{System Architecture}

\label{sec:architecture}

\subsection{Components}

The HMM system comprises six interacting components:

\begin{itemize}[leftmargin=1.1em]
    \item \textbf{Workers}: Provide compute capacity (GPU, CPU, RAM, bandwidth, storage). Each worker registers with a capability vector and stakes collateral.
    \item \textbf{Clients}: Request compute jobs, escrow payment tokens, and receive demand credits $\bm{p}$.
    \item \textbf{Routers}: Match jobs to resources via the convex path solver (\ref{eq:routing}).
    \item \textbf{HMM Pools}: Per-resource-class pools maintaining the Hamiltonian invariant.
    \item \textbf{Registry}: On-chain registry of job specifications, attestations, and settlements.
    \item \textbf{Validators}: PoAI verification nodes that sample-check job execution quality.
\end{itemize}

\subsection{Job Lifecycle}

\begin{enumerate}[leftmargin=1.4em]
    \item Client submits job specification with requirements $(\bm{r}, \bm{c})$ and locks collateral.
    \item Router queries HMM pools for a quote: $(\Delta\bm{q}, \bm{\pi}, f)$.
    \item Client accepts the quote; credits locked, resources allocated.
    \item Workers execute the job within a TEE and emit an attestation.
    \item Validators sample-check the attestation using PoAI protocol.
    \item Settlement: release payment to workers, rebate unused credits, distribute fees to LPs.
\end{enumerate}

\subsection{Dynamic Curvature Adaptation}

The curvature parameter $\lambda$ adapts to market conditions:
\begin{equation}
\lambda(t) = \lambda_0 \cdot \left(1 + \alpha \cdot \text{Vol}_{7d}(\Delta\bm{q}) + \beta \cdot \text{Util}(\bm{q})\right),
\end{equation}
where $\text{Vol}_{7d}$ is 7-day rolling volatility of trade volumes and $\text{Util}(\bm{q}) = 1 - \min_i(q_i / \bar{q}_i)$ is the utilization factor. High volatility or high utilization increases curvature, providing more stability at the cost of higher slippage.

